\section{Unweighted Setting}
\label{sec:unweighted}

In this scenario, we study the problem of packing a subset with a maximum number of hyperspheres from a collection $\mathcal{C}$ of unit weight $d$-dimensional hyperspheres in a unit bin. We call this problem unweighted $\mathsf{d-CPK}$ and its two-dimensional version unweighted $\mathsf{CPK}$. We propose a $5/2$-approximation algorithm for unweighted $\mathsf{d-CPK}$ that works for all dimensions $d$ based on ideas from the last section. Next, we introduce a new  $5/3$-approximation algorithm specifically for the two-dimensional case i.e. unweighted $\mathsf{CPK}$.

\subsection{$5/2$-approximation algorithm in all fixed dimensions $d$}
We now solve the $d$-dimensional hypersphere packing problem when each hypersphere has unit weight. To do this, we show that given a set of at most three hyperspheres that can be packed in a unit hypercube, we can always return the packing with a simple heuristic. %The proof of this is added to the appendix.z%
\begin{lemma}
	\label{twopack}
	Given a set $Q$ of at most two $d$-dimensional hyperspheres, there is a constant time algorithm to determine whether all hyperspheres of $Q$ can be packed inside a unit $d$-dimensional hypercube, and to return a packing if it exists. 
\end{lemma}
\begin{proof}
	Given an optimal packing of at most two hyperspheres, we can always push them to diagonally opposite corners of the unit hypercube.
	
	Our algorithm to find a packing of $Q$ is to simply to pack one hypersphere in a corner so that it touches all sides incident on that corner and then, if possible, pack the other hypersphere on the diagonally opposite corner so that it touches all sides incident on that corner. Return this packing if it was possible to pack $Q$ in this way, otherwise return that $Q$ cannot be packed inside a unit hypercube.
\end{proof}

\begin{theorem}
	There is polynomial time $5/2$-approximation algorithm for unweighted d-CSK in all dimensions $d$.
\end{theorem}
\begin{proof}
	We again get a packing in an augmented bin of profit of at least $w(\OPT)$ using the algorithm from lemma \ref{extended-wowa} for $\eps< 1/5^{2d+3}$. We then distinguish our algorithm between three parts depending on the value of $w(\OPT)$.
	
	Firstly, we can pack one hypersphere trivially and we can pack two hyperspheres using lemma \ref{twopack} for $2\leq w(\OPT)\leq 5$. 
	
	Next, when $6 \leq w(\OPT)$, we restrict $\eps<1/2d^2$ according to the bound in lemma \ref{uniqhuge} which guarantees us that for all $d$, there is only at most one hypersphere with a diameter greater than $1-\eps$. We then kill off this unique \textit{huge} hypersphere if it exists. We then apply the algorithm from lemma \ref{nohuge} as there is no \textit{huge} hypersphere and split our augmented packing into two packings into unit bins. We then return the bin with the larger weight. Note that we return at packing of at least $\lceil \frac{w(\OPT)-1}{2} \rceil$ hyperspheres. This always returns an output better than $2/5 \cdot w(\OPT)$ for $w(\OPT)\geq 6$. Combining the two cases, we get our result.
	
\end{proof}
\begin{comment}
	content...

\subsection{5/3-approximation algorithm for two dimensions}

In this section we attempt to get a 5/3-approximation algorithm for the unweighted $\mathsf{CSK}$ problem, i.e. the unweighted two-dimensional case of circle knapsack packing into a square. We begin by proving structural lemmas and then finally arrive at our algorithm.

\begin{lemma}
	\label{threepack}
	Given a set $Q$ of at most three $2$-dimensional circles, there is a constant time algorithm to determine whether all circles of $Q$ can be packed inside a unit square, and to return a packing if it exists. 
\end{lemma}
\begin{proof}
	%The idea is to first show that in any such packing, one circle can always be pushed to a corner. We can prove this combinatorially. Then the other two circles can be packed easily.%
	If the set \(Q\) contains only 1 circle, then if the radius of the circle is \(\leq \frac{1}{2}\), then we can always pack it with it's center overlapping with the center of the unit bin.
	
	If the set \(Q\) contains 2 circles, then we can pack the one circle at a corner touching two of the adjacent bin sides. Then we can pack the other circle at a diagonally opposite corner similarly touching two of the adjacent sides. If this strategy fails to produce a valid packing, then no such packing exists
	
	For the case, when the set \(Q\) contains 3 circles, then we show that all the possible valid packing can be modified to have a particular regular structure and we exhaust all of these possibilities to find out a feasible packing or conclude that no such packing exists.
	
	Suppose, the circles be \(Q=\{C_1,C_2,C_3\}\) and there is a valid packing named \(P\) of \(Q\) in a unit bin . We first consider a rectangular bin \(B\) whose each of the four sides touch at least one of the three circle. These is essentially the rectangle with the smallest side lengths which inscribe all the 3 circles. Firstly we note that there are only 3 circles and 4 sides of the rectangle which implies that at least one circle must touch two adjacent sides of \(B\). Without loss of generality we can assume that \(C_1\) is a circle which touch two adjacent sides of the bin \(B\). If there are more than 1 circle which touch two adjacent sides, then we consider the one with the smallest radius. Now, we consider the corner which has two of it's adjacent sides touched by \(C_1\) and name this corner point \(A^'\). We translate the bin \(B\) along with the circles in \(Q\) with their relative positions maintained inside the bin so that \(A^'\) overlaps with the it's nearest corner point \(A\) in the unit bin. After this step, we have \(C_1\) touching two of the adjacent sides of the unit bin. Depending on the relative ratio of the radius of \(C_2\) and \(C_3\) with \(C_1\), we can either translate each of \(C_2\) and \(C_3\) such that they either touch (i) two of the sides of the unit bin or (ii) one side of the unit bin and the circle \(C_1\). In both of the cases we have the position of the center of \(C_2\) and \(C_3\) deterministic depending upon which sides of the bin they touch. Now that we have shown this structural property for the feasible packing of set \(Q\), in order to decide if all the circles in \(Q\) can be packed in the unit bin we proceed as follows. For each \(i=1,2,3\) we pack the circle \(C_i\) in a corner \(A\) of the bin. For the other two circles, we try to pack them exhausting both of the cases described above for each of the circle. If any of the cases returns us a feasible packing we have found one feasible packing for \(Q\). Otherwise, we note that no feasible packing exist.
\end{proof}

%\label{greedyNFDR_section}

In the paper \cite{NFDH}, the authors consider the two-dimensional packing problem for packing squares into rectangles and propose a next fit decreasing height algorithm. We trivially extend their work for packing two dimensional circles with a bounded radius and unit weight into a unit square as \(\mathtt{NFDR}\) (Next fit decreasing radius algorithm), more formally stated as the lemma below.

\begin{lemma}
	\label{NFDR}
	Given a set of two-dimensional circles $Q$ such that each circle in $Q$ has a radius smaller than $\mu$, the
	\(\mathtt{NFDR}\) strategy packs an area of at least $\frac{\pi}{4}(1-4\mu)$ in a unit square.
\end{lemma}
\begin{proof}
	Draw a axis parallel bounding box over each of the circles we need to pack and pack this effective set of squares using the NFDH algorithm. In the paper \cite{NFDH} the NFDH algorithm produces the following guarantee, ``a set of squares $S$ with sidelengths $\leqslant \delta.$ If NFDH cannot place any square in a rectangle $R := r_1 \times r_2$, where $r_1, r_2 \geqslant\delta$, then total wasted space is less than or equal to $ \delta(r_1 + r_2)$.'' In our case, \(r_1=1,\) and  \(r_2= 1\) and \(\delta=2\mu.\) So, we get that total wasted space is less than or equal to \(4\mu\). Hence, we pack an area of \(1-4\mu\) at least. But, since we enclosed each circle with a square, we end up losing a fraction \(\frac{\pi}{4}\) of the area. Hence, finally we pack at least an area of \(\frac{\pi}{4}(1-4\mu).\)
\end{proof}

However, we generalise this for set of circles with unit weights with bounded radius $\leq \mu$. The appropriate modification to the problem is for us to have to pack a subset of circles $S'$ from a given set of circles $S$ with radii less than or equal to $\mu$ such that the number of circles that is packed is maximum. Let us call the following strategy $\mathtt{greedyNFDR}$:

We first sort all the circles in increasing order of their radius and declare a empty set $T$. We keep choosing the circles starting from the first circle and it to $T$ until we run out of circles or on adding a new circle to $T$, the sum of area of all the circles in $T$ exceeds $\frac{\pi}{4}(1-4\mu)$. Unless we run out of circles in constructing $T$, we can always construct $T$ with sum of area of circles in $T$ at least $\frac{\pi}{4}(1-4\mu)-\pi\mu^2$ since all the circles in $S$ have radius $\leq \mu$.

\begin{lemma}[{\(\mathtt{greedyNFDR}\)}]
	\label{greedyNFDR_lemma}
	The $\mathtt{greedyNFDR}$ strategy for circles with radii $\leq \mu$ and unit profit packs at least $\frac{\pi}{4}(1-4\mu)-\pi\mu^2$ fraction of circles that can be packed in the optimal packing of the square knapsack. 
\end{lemma}

\begin{proof}
	Let us order the circles in $S$ in increasing order of radius as \(S=\{c_1,c_2,c_3,...,c_n\}\). Suppose that the circle is able to pack a subset \(S_{\text{packed}}=\{c_1,c_2,...,c_s\}\) of size s where each circle has radius \(\leq r_{s}\) and is unable to pack the circles \(S_{\text{unpacked}}=\{c_{s+1},c_{s+2},...,c_n\}\) where each circles has radius \(\geq r_{s+1}\ \geq r_{s}\).  Note that if the total area of circles in $S$ is less than \(\frac{\pi}{4}(1-4\mu)-\pi\mu^2\), we can trivially pack all the circles in the bin using $\mathtt{greedyNFDR}$ and the theorem holds trivially since \(1 \geq  \frac{\pi}{4}(1-4\mu)-\pi\mu^2\) . Otherwise note that if $a_i$ denotes the area of the $i^{th}$ circle in $S$, \[\sum_{i=1}^{s}a_i \geq  \frac{\pi}{4}(1-4\mu)-\pi\mu^2\] Also since the total area of all the circles is less than 1 \[\sum_{i=s+1}^{n}a_i \leq  1-(\frac{\pi}{4}(1-4\mu)-\pi\mu^2)\]
	Also all the circles in \(S_{\text{packed}}\) have radius \(\leq r_{s}\) and among all the possible sets of \(S_{\text{packed}}\), $S$ will be minimum when all the circles in \(S_{\text{packed}}\) have radius \(r_{s}\). Hence we can bound s as \[s \geq \frac{\frac{\pi}{4}(1-4\mu)-\pi\mu^2}{\pi r_{\text{s}}^2}\]
	Also all the circles in \(S_{\text{unpacked}}\) have radius \(\geq r_{s}\) and among all the possible sets of \(S_{\text{unpacked}}\), the cardinality of the set \(S_{\text{unpacked}}\) will be maximum when all the circles in \(S_{\text{unpacked}}\) will have radius \(r_{s}\). Hence we can bound the cardinality of the set \(S_{\text{unpacked}}\) as \[n-s \leq \frac{1-(\frac{\pi}{4}(1-4\mu)-\pi\mu^2)}{\pi r_{\text{s}}^2}\] where $n$ is the total number of circles in $S$. Therefore the minimum value of \(\frac{s}{n}\) can be bound as
	\begin{align*}
		\frac{s}{n} \geq \frac{\frac{\frac{\pi}{4}(1-4\mu)-\pi\mu^2)}{\pi r_{\text{s}}^2}}{\frac{(\frac{\pi}{4}(1-4\mu)-\pi\mu^2)}{\pi r_{\text{s}}^2} + \frac{1-(\frac{\pi}{4}(1-4\mu)-\pi\mu^2)}{\pi r_{\text{s}}^2}} \\
		\Rightarrow \frac{s}{n} \geq \frac{\pi}{4}(1-4\mu)-\pi\mu^2
	\end{align*}
	Hence we will always be able to pack \(\frac{\pi}{4}(1-4\mu)-\pi\mu^2\) times the optimal total number of circles using $\mathtt{greedyNFDR}$.
\end{proof}







Let $\OPT$ be the optimal packing for $\mathsf{CSK}$ and $\mathcal{C}_{\OPT}$ be the set of circles in $\OPT$. 
We first partition set $\mathcal{C}$ of circles into three groups based on the size of their radii and an input parameter $\eps>0$ called small, medium, and large. We also assume \(\eps \leq 0.1\) and our algorithm successfully provides a valid packing for these bound on \(\eps\) We choose the bounds on radii for partitioning such that the medium circles have a total profit of at most $\eps\cdot w(\OPT)$ and hence ignore the medium circles from now on, losing a small amount of profit.

We then pack the large circles from $\mathcal{C}_{\OPT}$ in our unit bin using an approximation scheme and call the packing $P^+$. Depending on the number of large circles compared to number of small circles, we choose to pack different set of circles in our final packing. If there are not too many small circles in $\mathcal{C}$, we can simply pack them in the pocket created by three or more touching large circles in $P^+$ and pack the small circles inside the pocket using the algorithm mentioned in \cite{Fekete2019PackingDI} for packing disks into disks to further get packing  \(P^*\) from \(P^+\). If there are a large number of small circles, we can kill off all the large circles, losing only a small amount of profit and approximate pack only the small circles in \(\mathcal{C}\) into the unit bin to get our desired packing \(P^*\).
When there are only a small number of circles in $\mathcal{C}_{\OPT}$, we can simply pack using brute force using lemma \ref{threepack}. We prove that this scheme results in a worst case approximation ratio of $5/3$.
\subsubsection{Partitioning the set of circles} \label{partiotion section}
We partition our set $\mathcal{C}$ using ideas from section \ref{medium}. We select the function according to our needs, ensuring a significant gap in the radii of large and small circles. This choice is crucial for establishing various structural properties.\\
Define $\rhoi{0} = \frac{1}{2}$, and $\rhoi{k} = \frac{1}{2}\times\rhoi{k-1}^2\eps$ for all integers $k \geqslant 1$. Then,

\begin{equation}
	\label{rho_k_new}
	\rhoi{k} = \frac{\eps^{2^{k}-1}}{{2^{2^{k+1}-1}}}
\end{equation}

For each integer $k \geqslant 1$ we define the following sets of circles:
\begin{align*}
	S_k = \{c_i \in S \mid r_{i} \in (\rho_{k},\rho_{k-1}]\}
\end{align*}

 We then remove medium circles from further consideration and lose at most $\eps$ fraction of the total number of circles. In the next sections, we show how to modify $\OPT$ to produce a packing of circles in \(\mathcal{C}_{L}^*\)to a packing $P^+$ with a very regular structure. We will further extend \(P^+\) to a packing \(P^*\) with including \(\mathcal{C}_{L}*\) and \(\mathcal{C}_{S}^*\)

To describe the construction of $P^*$, we need to consider two separate cases: (i) when $\mathcal{C}_{L}^* = \varnothing$, and (ii) when $\mathcal{C}_{L}^* \neq \varnothing$
\subsubsection{Packing the Big Circles} \label{Big_Circle_Pack}

When $\mathcal{C}_{L}^* = \varnothing$,we first remove all the medium circles in $\OPT$ to get a set of small circles $\mathcal{C}_{S}^*$ of profit at least $(1-\eps)w(\OPT)$ in which all remaining circles have radius at most $\rho_{\tau} \leq \frac{\eps}{8}$, as $\tau \geq 1$. We will pack these circles using the greedyNFDR as described in lemma \ref{greedyNFDR_lemma}  for unit weight circles. Using this strategy we will be able to pack at least \(\frac{\pi}{4}(1-4\rho_{\tau})-\pi\rho_{\tau}^2\) times the total number of small circles in our optimal solution. Since \(\rho_{\tau} \leq \frac{\eps}{8}\).We will always be able to pack more than \(\frac{3}{5}\) times the total number of small circles in optimal solutionfor $\eps \leq 0.1$.

Now we deal the case when then number of circles in \(\mathcal{C}_{L}^*\) is non-zero. If the number of large circles in \(\OPT\) is less than 6 we proceed as follows. We first remove all the medium circles $\mathcal{C}_M ^*$ in \(\OPT\) from the packing. We also remove all the small circles $\mathcal{C}_S ^*$ temporarily for now. We will add them back later. Now we are only left with large circles in the bin. Then we can modify their packing using Lemma \ref{threepack} and pack 1 large circle when there is 1 large circle in \(\OPT\), 2 large circles when there is 2 large circles in \(\OPT\) or pack 3 large circles when there is either 3, 4, or 5 large circles in \(P\). 

\begin{figure}[ht]
	\begin{center}
			\includegraphics[width=6cm]{Images/4 by 3 apx figure.pdf}
		\caption{\(1\times 1\) bin with \(1\times\rho_{\tau-1}\) strips along the four sides to modify packing of large circles}
		\label{fig:4-by-3-apx-img}
		\end{center}

\end{figure}

Now we look at the case when, the number of large circles in \(\OPT\) is more than 5. Here is the scheme of things we follow to arrive at the structure we want. We begin with an optimal packing \(\OPT\) in the square bin \(ABCD\) as shown in figure \ref{fig:4-by-3-apx-img}. We first remove all the medium circles $\mathcal{C}_M ^*$ in \(\OPT\) from the packing. We also remove all the small circles temporarily for now. We will add them back later. Now we are only left with only circles in \(\mathcal{C}_{L} ^*\) in the bin. If the number of large circles is \(6\) or more circles.  We define the following cuts on the packing of \(\mathcal{C}_{L}^*\): \(\ell_1\) representing the line \(EG\), \(\ell_2\) representing the line \(GH\), \(\ell_3\) representing the line \(FH\), and \(\ell_4\) representing the line \(EF\). These cuts are made at a distance of \(\rho_{\tau-1}\) from the sides of the bin. Recall that all our large circles have a radius \(\geq \rho_{\tau-1}\). Also we assume that there exists no circle with radius \(\geq \frac{1}{2}-\rho_{\tau-1}\) that is no circle intersects two diametrically opposite strips. We will deal this case separately when there exists a circle with radius \(\geq \frac{1}{2}-\rho_{\tau-1}\). These cuts will generate four rectangular regions, \(R_1,R_2,R_3,\) and \(R_4.\) The rectangular regions include the shaded corners as well. The grey shaded corners are squares of dimension \(\rho_{\tau-1}\times\rho_{\tau-1}.\) These will be essential, as we prove important structural lemmas involving these four regions.

\begin{lemma}
	\label{only4-circles}
	At most one large circle can intersect each of the grey regions. In other words, at most 4 large circles can intersect the regions shaded grey. 
\end{lemma}
\begin{proof}
	To prove the claim, we will analyze the grey region on bottom left between point C,G and same arguments can be used for the other grey squares as well.
	For a large circle to intersect the grey squares, its center must lie in a `feasible region'. For a (large) circle of radius \(\rho_{\tau-1}\) and the region enclosed between point \(G,C\) the feasible region would be a quarter circle centered at \(G\) of radius \(\rho_{\tau-1}\) since if the center lies beyond a distance of \(\rho_{\tau-1}\) from \(G\) then, it's not possible for it to intersect the grey square. 
	
	Consider the feasible regions for the centers of two circles of radius \(r_1,r_2\geqslant\rho_{\tau-1}\), the feasible region would be quarter of a circle. For example, consider the circle with radius \(r_1\) and focus on it's center. One case where the circle intersects the maximum of the grey square is when it's center lies at the point \((r_1,r_1)\). The other extreme case will be when the circle intersects just \(G\) as in Fig. \ref{fig:4-by-3-apx-img}. And rest of the conditions will be imposed by the fact that the circle cannot go out of the bin. So, this will end up tracing a quarter circle where the arc will be defined by the equation \((x-\rho_{\tau-1})^2+(y-\rho_{\tau-1})^2=r_1^2\), and the lines \(x=r_1\), \(y=r_1\). The intersection of these will trace out the quarter circle. 
	
	Let the feasible region of the center of the circle with radius \(r_1\) be \(FR_1\), for \(r_2\) be \(FR_2\). Observe that the center of the feasible regions will always lie on the line \(y=x\) because of the above construction. To have a non-intersecting configuration of circles there needs to be at least a distance of \(r_1+r_2\) between their centers. When \(FR_1\cap FR_2 = \varnothing\), we claim that the maximum distance between any two points in \(FR_1\) and \(FR_2\) is less than the required \(r_1+r_2\) and thus, we cannot have a non intersecting configuration of two circles with both of them intersecting the grey square. From preliminary observation we can see that the two points furthest apart from each other must lie on the part defined by straight lines for \(FR_1\) and the arc for \(FR_2.\) 
	
	\begin{figure}[ht]
		\centering
		\includegraphics{Images/large_circle_corner_lemma.pdf}
		\caption{(a) The orange region highlights the feasible region for the dashed circle. (b) An illustrative example of the points and lines required for the proof of claim.}
		\label{fig:enter-label}
	\end{figure}
	\begin{claim}
		The center of \(FR_1\) and the point where \(y=x\) intersects \(FR_2\) are the two points which are furthest apart from each other.
	\end{claim}
	\begin{claimproof}
		Let, \(X\) be the center of \(FR_1\) and $Y$ be the point where \(y=x\) intersects the arc of \(FR_2.\) Now consider any point \(X'\) either on \(x=r_1\) or \(y=r_1\) on \(FR_1\) and \(Y'\) on the arc of \(FR_2.\) The shortest distance between both of them will be the straight line joining them. Now, join \(X\) and \(Y\) along a straight line and let it intersect \(FR_1\) at \(P\) and \(FR_2\) at \(Q\), now observe that, 
		\begin{align*}
		X'Y' &\leqslant X'P+PQ+Y'Q \\
		&\leqslant XP+PQ+YQ \\
		&= XY.
		\end{align*}
		Thus giving us that the longest distance between any two points on \(FR_1\) and \(FR_2\) when \(FR_1\cap FR_2=\varnothing\) is \(XY.\) 
	\end{claimproof}
	Now using standard cartesian coordinate system, we can figure out that the coordinates of \(X\) are \((r_1,r_1)\) and the coordinates of \(Y\) are \((r_2/\sqrt{2}+\rho_{\tau-1},r_2/\sqrt{2}+\rho_{\tau-1})\) due to the construction of \(FR_2.\) Now, applying the distance formula we prove that \(XY<r_1+r_2\) always. Let \(d\) be the distance/length of \(XY\), then
	
	\begin{align*}
	d^2 &= \Le(\frac{r_2}{\sqrt{2}}+\rho_{\tau-1}-r_1\Ri)^2+\Le(\frac{r_2}{\sqrt{2}}+\rho_{\tau-1}-r_1\Ri)^2 \\
	&\approx \Le(\frac{r_2}{\sqrt{2}}-r_1\Ri)^2+\Le(\frac{r_2}{\sqrt{2}}-r_1\Ri)^2 \\
	&= 2\Le(r_1^2+\frac{r_2^2}{2}-r_1r_2\sqrt{2}\Ri) = r_1^2 + (r_1^2+r_2^2-2\sqrt{2}r_1r_2)\\
	&\leqslant r_1^2 + (r_1^2+r_2^2 - 2r_1r_2) \\
	&= r_1^2 + (r_2-r_1)^2 \\
	&\leqslant r_1^2+r_2^2 \\
	&< r_1^2+r_2^2+2r_1r_2 = (r_1+r_2)^2
	\end{align*}{\tiny }
	Thus, we have that \(d\leqslant(r_1+r_2).\) Now the case then \(FR_1\subset FR_2\) is trivial to observe. Since one feasible region is contained in another we cannot have any selection of two points such that their distance is greater than \(r_1+r_2.\) Similarly for the case when \(FR_1\cap FR_2 \neq \varnothing\) observe that the the distance between the previously defined points \(X\) and \(Y\) decrease because of the two feasible regions intersecting each other. So naturally, the length of \(XY\) decreases as well, and since in the previous case we saw that \(d< r_1+r_2,\) this will apply to this case as well. Hence, we have that at most one large circle can intersect the grey squares.
\end{proof}
We call the circles which intersect at least two strips, or equivalently the grey region as shown in figure \ref{fig:4-by-3-apx-img}, \textit{corner circles}.
\begin{lemma}
	\label{delete-one-repack}
	{If there are many circles along with the corner circles which intersect two adjacent strips (\(R_i\) and \(R_{i-1}\) or \(R_{i+1}\)) we can always remove the largest circle in \(\mathcal{C}_{L}^*\) which intersect the strips and pack the circles which intersect two strips simultaneously into the space occupied by the removed circle.}
\end{lemma}
\begin{proof}
	Without loss of generality, we will consider the case of the strips \(R_1\) and \(R_4\). In order to prove this lemma, first let us consider the ordering of the cirlces which intersect two adjacent strips \(R_1\) and \(R_4\) as \(\mathcal{C}_{\text{intersect}}=\{C_1,...,C_m\}\) in the increasing order of their corresponding centers' distance from the point \(A\). Let us also denote the radius of circle \(c_i\) with \(r_i\). In order to prove our lemma we will find the lower bound for \(\frac{r_i}{r_{i-1}}\). First we observe that all the circles in \(\mathcal{C}_{\text{intersect}}\) have radius \(\geq \rho_{\tau-1}\). For a particular fixed circle \(C_i\) with radius \(r_i\), the radius of the \(i-1\)-th circle \(C_{i-1}\) will be maximum when the \(i\)-th circle and \(i-1\)-th circle touch two opposite boundaries of the bin and only touch the other side strip as shown in Figure \ref{fig:enter-label}.
	\begin{figure}[H]
		\centering
		\includegraphics[width=0.5\textwidth]{Images/consequitive_circle_in_square} % Scale image to the width of the text
		\caption{Two consecutive circles intersecting two adjacent strips R1 and R4}
		\label{fig:enter-label}
	\end{figure}
	Notice that in this case, we have \(O_nO_{n-1}=r_i+r_{i-1}\), \(O_nE=r_i-r_{i-1}-\rho_{\tau-1}\) and \(O_{n-1}E=r_i\). Applying Pythagoras theorem on triangle \(O_nO_{n-1}E\), we have

\begin{gather*}
	(r_i+r_{i-1})^2 = (r_i-r_{i-1}-\rho_{\tau-1})^2 + (r_i)^2 \\
	r_i^2 - 2\rho_{\tau-1}r_i + 2\rho_{\tau-1}r_{i-1} - 4r_ir_{i-1} + \rho_{\tau-1}^2 = 0
\end{gather*}
 For the ratio \(\frac{r_i}{r_{i-1}}\), consider two boundary cases. First, we consider the case when \(r_{i-1}=\rho_{\tau-1}\), that is \(r_{i-1}\) reaches it's minimum possible value. In this case we get, \(r_i=(3+\sqrt{6})\rho_{\tau-1}\). In this case \(\frac{r_i}{r_{i-1}}=3+\sqrt{6}>4\). Second we consider another boundary case, that is \(r_i=\frac{1}{2}\). In this case \(r_{i-1}=\frac{(\frac{1}{2}-\eps)^2}{2(1-\eps)}\). So in this case, we have \(\frac{r_i}{r_{i-1}}=\frac{(1-\eps)}{(\frac{1}{2}-\eps)^2}\geq 4\). It follows that \(\frac{r_i}{r_{i-1}}\geq 4\),we first note that \(r_i \leq r_n (\frac{1}{4})^{m-i}\) and we can make the following claim about the sum of area of circles from \(1 \text{ to } m-1\) in comparison to the area of the \(m\)-th circle. Denote the area of the \(i\)-th circle by \(A_i\). It follows that
 
\begin{center}
	\(\sum_{i=1}^{m-1} A_i \leq \sum_{i=1}^{m-1} \pi \left(r_m\left(\frac{1}{4}\right)^{m-i}\right)^2 = \sum_{i=1}^{m-1} \pi r_m^2 \left(\frac{1}{16}\right)^i < \frac{\pi r_m^2}{15} = \frac{A_m}{15}\)
\end{center}
	Now consider an inscribed axis parallel square inside the circle with radius $r_m$. This has an area of $2 r_m^2$. Using the Splitpacking algorithm from \cite{10.1007/978-3-540-39658-1_20}, we know that since the total area of the circles with radii $r_1,..., r_{m-1}$ is at most $0.53\cdot2r_m^2$, we can pack them inside the inscribed square in polynomial time.
\end{proof}
Now let us handle the case when there exists one huge circle in \(\mathcal{C}_{L}^*\) with a diameter greater than or equal to \(1-2\rhoi{\tau-1}\). Getting a partition in this case using our cuts is difficult because the \textit{huge} circle will intersect all of the \(R_i\)'s. But, in this case we can just delete the large circle from our packing \(P\) altogether. Notice that the area remaining after deleting the huge circle is atmost \(1-\frac{\pi}{4}(1-2\rhoi{\tau-1})^2\). Since \(\tau \geq 2\), we have \(\rhoi{\tau-1}\leq\frac{\eps}{8}\)\(\) and so we can bound the area of the remaining cirlces as \(1-\frac{\pi}{4}(1-2\rhoi{\tau-1})^2 \leq 1-\frac{\pi}{4}(1-\frac{2\eps}{8})^2 \leqslant 0.53 \). Now, using \cite{SplitPack} which claims that ``any set of circles with a combined area of up to 53.90\% of the square’s area can be packed" we can generate the approximate packing \(P^*\) which our algorithm will return. This \(P^*\) will have a profit of \(w(\OPT)-1.\) 

Now we consider the case when there is no large circle with diameter greater than or equal to \(1-2\rhoi{\tau-1}\).  Recall the rectangular strips \(R_i\)'s generated before. We now define \(C_1,C_2,C_3\) and \(C_4\) to be the set of disks in \(\mathcal{C}_{L}^*\) which only intersect only one of the above rectangular strips. If a circle intersects only one of \(R_i\)'s then it's included in the set \(C_{i}\). If a disk intersects both \(R_1\) and \(R_2\) then it's not a part of \(C_1\cup C_2\cup C_3\cup C_4.\) We will give a scheme to get rid of all such disks by removing at most one circle from each of the corner. The above creates a partition which will help in the analysis. Put the disks not in any of the \(C_i\)'s above in a set \(C_X.\) 

We do the following modification to the packing of large circles
\begin{enumerate}
	\item If we observe multiple disks which intersect or touch two adjacent rectangular strips then we use Lemma \ref{delete-one-repack} to handle this case and obtain the complete partition by removing the largest of such intersecting disks. The circles now do not intersect two adjacent strips anymore and now they are contained in \(C_X\). Notice that in this case, we get rid of the problematic case of circles intersecting multiple strips by removing at most one circle from each of the corner.\label{step3}
	\item Otherwise if we observe that if there is only a single disk which intersect or touch two adjacent rectangular strips or the shaded grey region then we just remove the circle and obtain the complete partition. In this case also we get rid of the problematic case of circles intersecting multiple strips by removing at most one circle from each of the corner.\label{step4}
	\item Partition the disks into sets \(C_1,C_2,C_3,C_4,\) and \(C_X.\) Where each set \(C_i\) contains those disks which only intersect \(R_i.\) \(C_X\) contains those which don't belong to any of the other \(C_i\)'s.
	\item Among \(C_1,C_2,C_3,C_4\) we remove the `cheapest' set of disks and also remove the \(1\times\rhoi{\tau-1}\)-size strip corresponding to that set, which in this case means removing the set with the least number of disks. We get a packing of the other 3 sets along with the set \(C_X\) in a bin of size \(1\times(1-\rhoi{\tau-1})\). Now we add back the circles removed in step \ref{step3} and set \ref{step4} to at least two of the corner of the \(1\times(1-\rhoi{\tau-1})\) sized bin.
\end{enumerate}

\begin{lemma}
	\label{rationalization-1-eps}
	We can rationalize the large circles' coordinates in the cuts generated \(1\times(1-\rhoi{\tau-1})\) bin by augmenting the bin by an \(\rhoi{\tau-1}\) factor. That is, we can augment the circles to have rational centers in a \(1\times1\) bin.
\end{lemma}
\begin{proof}
	We call the  \(1\times(1-\rhoi{\tau-1})\) bin as R.
	Let us begin by bounding \(n\), i.e., the number of large circles possible. To do this we make the observation that the total area of the large circles cannot exceed the area of \(R\). So we have \(n\pi(\rhoi{\tau-1})^2<1-\rhoi{\tau-1}\). This gives us, \(n<({1-\rhoi{\tau-1}})/{\pi\rhoi{\tau-1}^2}.\)  
	
	Now, we focus on rationalizing the \(x\)-coordinates. Number the circles in increasing value of $x$-coordinate from left to right. If two circles have the same $x$-coordinate, then order the one with the lower $y$-coordinate before the one with the higher $y$-coordinate. Let the $x$-coordinates of the $n$ circles in $P$ be $x_1,x_2,\ldots,x_n$. We propose the following (re-)packing scheme. For circle \(i\in P\) shift it to the right by \(\Delta x_i\) amount such that the shifts themselves follow the order \(\Delta x_1 < \Delta x_2 < \cdots <\Delta x_n\). Observe that this shifting scheme preserves the non-intersecting nature of the packing. 
	
	Superpose a semi-grid of side length \((\rhoi{\tau-1})^k\) on the bin, i.e., we superpose a grid (in which distance between two vertical lines is \((\rhoi{\tau-1})^k\) on the bin we have) where $k$ is a given integer. Observe that any line segment of length \((\rhoi{\tau-1})^k\) will have at least one integral multiple of \((\rhoi{\tau-1})^k\),  and hence we can get a suitable \(\Delta x_1 < \rhoi{\tau-1}^k\) such that the first circles aligns itself up on a grid line. We can similarly also find a \(\Delta x_2 < (\Delta x_1 + \rhoi{\tau-1}^k)+2\rhoi{\tau-1}^k\) where the \(2\rhoi{\tau-1}^k\) term is added to be used later as extra space. Note that the extra space we added leads to the property that $x_i- x_{i-1} \geqslant 2\rhoi{\tau-1}^k$ for $i=2,3,...,n$. Extending this for the \(n^{th}\) circle, we get \(\Delta x_i<(3i-2)\rhoi{\tau-1}^k.\) The only constraint is that the total displacement of the circles at the end should be at most \(\rhoi{\tau-1}\), so we get the following inequality as a sufficient condition for us to find a integral $x$-coordinate for every circle using the shifting scheme described above-  
	\[\Delta x_n<\rhoi{\tau-1}\Rightarrow (3n-2)(\rhoi{\tau-1})^k<\rhoi{\tau-1} \Rightarrow n < \frac{2}{3}+\frac{1}{3(\rhoi{\tau-1})^{k-1}}.\] Notice that this condition is satisfied from the area bound on $n$ at the beginning of the section for all $k \geq 3$ and a $\rhoi{\tau-1} < 0.5$, which is guaranteed as we have a bound on $\eps$ later on in the paper. Hence we can always shift the x-coordinates of the circles to fit themselves on a $(\rhoi{\tau-1})^k$ grid for $k \geq 3$.
	
	Intuitively, moving circles to the right in the $x$-direction will free up some space in the $y$-direction for the circles to move. Let us assume that this space from the $x$-movement above allows us to fit everything to an $\rhoi{\tau-1}^{\ell}$-grid in the $y$-direction (i.e., the grid-lines are horizontal now) with $\ell\geqslant k$ from above. We are now superposing a finer grid on the bin than the one we superposed above. 
	
	\begin{figure}[ht!]
		\centering
		\includegraphics[width=7cm]{Images/rational.png}
		\caption{The \(i^{\text{th}}\) circle and it's neighbors.}
		\label{fig:circles_rational}
	\end{figure}
	
	Now, consider the vertical movement to find integral coordinates for the $i^{\text{th}}$ circle.  Assume that it has $N$ ``neighbors" in the packing $P$. Neighbors are defined as the circles with an index larger than $i$ that disallow movement of the $i^{\text{th}}$ circle in the vertical $y$ direction by $\rhoi{\tau-1}^{\ell}$ (see Fig. \ref{fig:circles_rational}). The $i^{th}$ circle and its neighbors have the $y$-coordinates $y_i$ and $Y_1, Y_2, \cdots Y_N$. Let our rearrangement scheme move each of the circles by $\Delta y_i$ and $\Delta Y_1, \Delta Y_2, \cdots, \Delta Y_N$, respectively.
	The $i^{\text{th}}$ circle can have a maximum radius of $\frac{1}{2}$ as it is packed inside a $1\times 1$ square. The $x$-coordinate of the center of its $N$ neighbors have to be to the right of the $x$-coordinate of the $i^{\text{th}}$ circle. Because these are the only circles whose movement can free up space for the \(i^{\text{th}}\) circle to move in the \(y\)-direction. Since each of the $N$ neighbors have a radius of at least $\rhoi{\tau-1}$, we get that $N<\frac{2}{\rhoi{\tau-1}}$ from perimeter equations for $\rhoi{\tau-1}<0.5$, which is guaranteed by the bound on $\eps$ later.
	
	Due to the additional `\(2\rhoi{\tau-1}^k\)' factor which we added above in the rationalising scheme for the $x$-coordinates, we can claim that for each neighbor $j$ of the \(i^{\text{th}}\) circle we have $X_j- x_i \geq 2 (\rhoi{\tau-1})^k$. Since the circles were tightly packed before and now have an extra space of $X_j- x_i \geq 2 (\rhoi{\tau-1})^k$, using the expansion of cosine for small angles, we get that
	% \textcolor{red}{Using simple trigonometric arguments} \textcolor{blue}{[P: please add and and all arguments and calculations explicitly]}, this leads to 
	the space for movement in the y-direction is  such that $|\Delta Y_j|+ |\Delta y_i| \geq (\rhoi{\tau-1})^{2k}$.
		
	% After the circle has moved along the \(x\)-direction, recall the fact that the shortest distance between two circles is \(C_1C_2-r_1-r_2\), where \(C_1C_2\) is the distance between the centers of the two circles and \(r_1,r_2\) are their respective radii. So, after joining both the centers using a line segment observe that the `allowed' \(y\)-axis movement will be the vertical projection of the ``shortest distance" as highlighted by the small black triangle in Figure \ref{fig:tan_argument}. Consider the triangle formed by the old center \((C_1^o)\), the new center \(C_2^n\) (of the small circle) and the center of the big circle \((C)\). Observe that, \(C_1^oC=r_1+r_2\) where \(r_1\) is the radius of the large circle and \(r_2\) is the radius of the small circle and \(C_1^oC_1^n\geqslant 2\rhoi{\tau-1}^k\). Using this information, and applying the sine and cosine rules repeatedly, we can figure out the angle \(\angle C_1^oCC_1^n\) and length \(CC_1^n\). Using that, we can calculate the vertical projection of the smallest distance between the circles, which comes out to be greater than \(2\rhoi{\tau-1}^{2k}.\) 
	Now, we have to consider the vertical movement of each of the \(N\) neighbors. If $l \geq 2k+1$, then we have at least $\frac{2}{\rhoi{\tau-1}}$ integral points in the vertical free space ($2(\rhoi{\tau-1})^k$) for the $i^{th}$ circle. So we can find a $\Delta y_i$ such that $y_i$ lands on an integral coordinate. We do this for each circle $i$. In addition, we are guaranteed that for movements in the $N$ neighbors in further steps, we are sure to find a free integral y-coordinate as $N<\frac{2}{\rhoi{\tau-1}}$.
	
	% Since we have generated a `vertical free space' of \(2\rhoi{\tau-1}^{2k}\), so in a \(\rhoi{\tau-1}^{2k+1}\) grid, we will have \(\frac{2\rhoi{\tau-1}^{2k}}{\rhoi{\tau-1}^{2k+1}}=\frac{2}{\rhoi{\tau-1}}\) integral grid points. Now, we know \(N<\frac{2}{\rhoi{\tau-1}}\), so we can find at least one integral grid point for the \(N+1\) (`+1' because we add the `\(i^{\text{th}}\)' circle itself as well) circles under consideration. Hence, we observe that any grid of width \(\ell\geqslant 2k+1\) will work.
	Now, for concrete values we substitute \(k=3\) which gives us that everything can be fit into a \(\rhoi{\tau-1}^{2\cdot 3+1} grid.\) These values can be optimised a bit but we do not care too much about that. We can see that everything can be fit into an $\rhoi{\tau-1}^7$-grid after augmenting one of the sides of the square.
\end{proof}
% \subsection{Packing the big circles onto the grid}
%Now after rationalisation, we pack the big circles onto the grid in constant time. To prove that we can do so in constant time, we first note that the number of big circles is bounded by a constant as each big circle has a radius of at least $\eps$ and the total area of the bin is $1$. Now, there are again a constant number of subsets possible of the set of all big circles. Using rationalisation in coordinates in the augmented bin, we achieve constant number of possible packings for each subset of the big circle. Henceforth, we can find a rationalised packing of the big circles with maximum profit in constant time and proceed with it for the rest of the procedure. %
We had a packing of a subset of \(\mathcal{C}_{L}^*\) in a bin of size  \(1\times(1-\rho_{\tau-1})\). We will now augment the bin to have dimension \(1\times1\) with all the circles having rational coordinates to their corresponding centers using lemma \ref{rationalization-1-eps}. We call this packing \(P^+\).
The final structural claim we make about the weight of \(P^+\) is that the modified packing is guaranteed to have weight 
	\[
	w(P^+)= 
	\begin{cases}
	\lvert \mathcal{C}_{L}^* \rvert, & \lvert \mathcal{C}_{L}^* \rvert\leq 3 \\
	3, & 4 \leq \lvert \mathcal{C}_{L}^* \rvert \leq 5 \\
	\lceil\frac{3}{4}\cdot\lvert \mathcal{C}_{L}^* \rvert\rceil-1 & 6 \leq \lvert \mathcal{C}_{L}^* \rvert
	\end{cases}
	\]
    In the case when the number of large circles is $<6$, we use lemma \ref{threepack} , and the profit of the packing is straightforward from the packing strategy. In the other case, we note that,
	there can be a maximum of 4 circles removed from each of the corners during generating the partitioning, we are left with at least \(\lvert \mathcal{C}_{L}^* \rvert-4\) disks. Now, we remove the set \(C_i\) \(i\in\{1,2,3,4\}\) which has the least number of disks (which translates to the cheapest set among all \(C_i\)'s). This will result in loss of atmost \(\left\lfloor\frac{\lvert \mathcal{C}_{L}^* \rvert-4}{4}\right\rfloor\) profit. After the removal we add atleast two of the circles which we have earlier removed during the partitioning in their original positions and so the final number of large circles in \(P^+\) is atleast \[\lvert \mathcal{C}_{L}^* \rvert-4+2-\left\lfloor\frac{\lvert \mathcal{C}_{L}^* \rvert-4}{4}\right\rfloor = \left\lceil\frac{3}{4}\cdot\lvert \mathcal{C}_{L}^* \rvert\right\rceil-1\]
\subsubsection{Packing the small circles} \label{Small-Circle-Pack}
Now we will add back the small circles in \(P^+\) to get the final packing structure \(P^*\). Let us denote the total area of small circles in \(\mathcal{C}_{S}^*\) as \(A_{\text{small}}\) and the total area of large circles in \(\mathcal{C}_{L}^*\) as \(A_{\text{large}}\).
% Consider two separate cases as follows.

\textbf{Case - I : } \(A_{\text{small}} \leq A_{\text{large}}\times \frac{\rho_{\tau-1}^2}{30}\)

When the combined area of small circles in \(\mathcal{C}_{S}^*\) is very small as compared to the total area of large circles, we can pack all the small circles in the pocket created by three or more touching large circles in \(\mathcal{C}_{L}^*\) or pack them by removing at most one large circle.

First, consider the case when the number of large circles in \(\mathcal{C}_{L}^*\) is 1. In this case, the large circle will have a radius \(\leq \frac{1}{2}\). If our packing has a large circle of radius at most \(\frac{1}{2}\), then we can have a imaginary circle of radius at least\(\frac{\sqrt{2}-1}{2(\sqrt{2}+1)}\) touching two sides of the bin at one of the corners and not intersecting the large circle. We know that the total area of small circles is at most \(\pi (\frac{1}{2})^2 \cdot \dfrac{\frac{1}{4}}{30}=\frac{\pi}{480}\). The imaginary corner circle has area at least \(\pi (\frac{\sqrt{2}-1}{2(\sqrt{2}+1)})^2\). Using the algorithm in \cite{Fekete2019PackingDI} for packing circles into circles, we can pack all the small circles in the imaginary circles since \(\frac{\pi}{480}\leq \frac{1}{2}\cdot\pi(\frac{\sqrt{2}-1}{2(\sqrt{2}+1)})^2\) and we get packing for all the small circles in \(\mathcal{C}_{S}^*\) along with the large circle in \(\mathcal{C}_{L}^*\).

For the case when number of large circles in \(\mathcal{C}_{L}^*\) is 2. We pack the small circles using different approaches depending upon the number of small circles. If the number of small circles in \(\leq 3\). We have at most 5 circles and we can pack them in the bin using lemma \ref{threepack} and always pack 3 out of at most 5 circles.

Other wise in this scenario we choose the biggest of the two large circles that is already packed in our bin. Note that this circle has radius at least \(\rho_{\tau-1}\). Consider the square of maximum area that can be inscribed this circle and has sides parallel to that of the bin. We replace the circle with this square and name it as \(S\). This square has side length at least \(\rho_{\tau-1} \times \sqrt{2}\). Also note that \(A_{\text{large}} \leq 1\)  implies \(A_{\text{small}} \leq \frac{\rho_{\tau-1}^2}{4}\) and all the small circles have radius \(\leq \rho_{\tau}\). We pack all the small circles in \(S\) using NFDR. For circles with radius less than \(\rho_{\tau}\), we can pack either all the circles or circles with total area at least \(\frac{\pi}{4}(2\rho_{\tau-1}^2-4\sqrt{2}\rho_{\tau-1}\rho_{\tau})\)
in a square with side length \(\rho_{\tau-1}\times\sqrt{2}\). Note that \(\rho_{\tau} = \frac{1}{2} \times \rho_{\tau-1}^2\eps\). 
Hence we can see
\begin{align*}
\frac{\pi}{4}(2\rho_{\tau-1}^2-4\sqrt{2}\rho_{\tau-1}\rho_{\tau})=\frac{\pi}{4}(2\rho_{\tau-1}^2-2\sqrt{2}\rho_{\tau-1}^3\eps) \geq \frac{\rho_{\tau-1}^2}{30}
\end{align*}
 Hence, we can pack all the small circle in \(P\) by replacing one big circle from our bin.

For the case when the number of large circles in \(\mathcal{C}_{L}^*\) is \(\geq 3\), we will pack all the small circles in \(\mathcal{C}_{S}^*\) in pocket created by three or more touching large circles. In the packing \(P^+\) of large circles in \(\mathcal{C}_{L}^*\), all the circles have radius \(\geq \rho_{\tau-1}\). In this packing of the large
we can find a pocket created by three or more touching large circles and embed an imaginary circle in it . In order to provide a feasible packing of small circles, we need to find lower bound for the radius of the imaginary circle in the pocket. We can note that the imaginary circle in the pocket will have the least radius when it is created by three mutually touching circles of radius \(\rho_{\tau-1}\).
%\begin{center}
%	\begin{tikzpicture}
%	% Define the radius of the large circles
%	\def\Radius{2cm}
%	% Calculate the radius of the small circle
%	\def\smallRadius{(\Radius*(2-sqrt(3))) / sqrt(3)}
%	
%	% Define coordinates for the centers of the large circles
%	\coordinate (X) at (0,0);
%	\coordinate (Y) at ($(X) + (60:2*\Radius)$);
%	\coordinate (Z) at ($(X) + (0:2*\Radius)$);
%	
%	% Draw the large circles and label their centers
%	\draw (X) circle (\Radius);
%	\draw (Y) circle (\Radius);
%	\draw (Z) circle (\Radius);
%	\node at (X) [below] {X};
%	\node at (Y) [above] {Y};
%	\node at (Z) [below] {Z};
%	
%	% Calculate the center of the small circle (centroid of the triangle)
%	\coordinate (O) at (barycentric cs:X=1,Y=1,Z=1);
%	
%	% Draw the small circle and label its center
%	\draw (O) circle (\smallRadius);
%	\node at (O) [below] {O};
%	\end{tikzpicture}
%\end{center}
\begin{figure}[ht]
	\begin{center}
		\includegraphics[width=7cm]{Images/3-circles-touching.pdf}
		\caption{Packing small circles in the \textit{pockets} of the large circle packing}
		\label{fig:pocket}
	\end{center}
	
\end{figure}
If the circles with center at \(X,Y,Z\) each have radius \(\rho_{\tau-1}\). Then the circle with center O in the pocket created by the three big circles will have radius \(\rho_{\tau-1}\times\frac{2-\sqrt{3}}{\sqrt{3}}\).
This is the least possible radius of a circle in a pocket created by three or more circles. Note that in this case the total area of small circles is \(A_{\text{small}} \leq A_{\text{large}}\times \frac{\rho_{\tau-1}^2}{30}\leq\frac{\rho_{\tau-1}^2}{30}\). Using the algorithm in \cite{Fekete2019PackingDI} for packing disks into disks, we can pack all the small circles in the imaginary disk in the pocket since \(\frac{\rho_{\tau-1}^2}{30}\leq \frac{1}{2}\times\pi (\rho_{\tau-1}\times\frac{2-\sqrt{3}}{\sqrt{3}})^2\) and we get packing for all the small circles in \(\mathcal{C}_{S}^*\) along with the existing packing a subset of large circles \(\mathcal{C}_{L}^*\) in \(P^+\).

In this case we analyze the weight of final packing \(P^*\) in comparison of \(P\) in different cases: 

\[
w(P^*) = 
\begin{cases}
\lvert \mathcal{C}_{L}^* \rvert+\lvert \mathcal{C}_{S}^* \rvert, &  \lvert \mathcal{C}_{L}^* \rvert=1 \\
\lvert \mathcal{C}_{L}^* \rvert+\lvert \mathcal{C}_{S}^* \rvert, & \lvert \mathcal{C}_{L}^* \rvert=2 \text{ and } \lvert \mathcal{C}_{S}^* \rvert= 1
\\
3, &\lvert \mathcal{C}_{L}^* \rvert=2 \text{ and } \lvert \mathcal{C}_{S}^* \rvert= 2 \text{ or } 3 \\
\lvert \mathcal{C}_{L}^* \rvert-1 + \lvert \mathcal{C}_{S}^* \rvert,& \lvert \mathcal{C}_{L}^* \rvert=2 \text{ and } \lvert \mathcal{C}_{S}^* \rvert\geq 4 \\
3+\lvert \mathcal{C}_{S}^* \rvert, & 3 \leq \lvert \mathcal{C}_{L}^* \rvert \leq 5
\\
\lceil\frac{3}{4}\cdot\lvert \mathcal{C}_{L}^* \rvert\rceil-1+\lvert \mathcal{C}_{S}^* \rvert & 6 \leq \lvert \mathcal{C}_{L}^* \rvert

\end{cases}
\]

In this case our packing \(P^*\) has profit \(\geq \frac{3}{5}\OPT\) for all the subcases

\textbf{Case - II : } \(A_{\text{small}} > A_{\text{large}}\times \frac{\rho_{\tau-1}^2}{30}\)

When the number of small circles in \(P\) is much more compared to number of big circles in \(P\), we can get large profit by packing only the small circles even after discarding the large circles. We will use this idea in this scenario. Let's denote the number of small circles and large circles in \(P\) as \(n_{\text{small}}\) and \(n_{\text{large}}\). We have \(\sum_{c_{i} \in \mathcal{C}_{S}^*} a_{i} \geq (\sum_{c_{i} \in \mathcal{C}_{L}^*} a_{i})\times \frac{\rho_{\tau-1}^2}{30}\). We try to find the bound on minimum value of the ratio \(\frac{n_{\text{small}}}{n_{\text{large}}}\). Denote the circles in \(\mathcal{C}_{S}^*\) as \(\mathcal{C}_{S}^*=\{C_{s_{1}},C_{s_{2}},C_{s_{3}},...,C_{s_{n_{\text{small}}}}\}\) and circles in \(\mathcal{C}_{L}^*\) as \(\mathcal{C}_{L}^*=\{C_{l_{1}},C_{l_{2}},C_{l_{3}},...,C_{l_{n_{\text{large}}}}\}\) 

Note that \[A_{\text{small}}=\sum_{i=1}^{n_{\text{small}}} a_{s_{i}}\] 
and all the circles in \(\mathcal{C}_{S}^*\) have radius \(\leq \rho_{\tau}\). For a fixed \(A_{\text{small}}\) and among all the possible candidate sets of \(\mathcal{C}_{S}^*\), \(n_{\text{small}}\) will be minimum when all the circles in \(\mathcal{C}_{S}^*\) have radius \(\rho_{\tau}\). Hence we can bound \(n_{\text{small}}\) as \[n_{\text{small}} \geq \frac{A_{\text{small}}}{\pi \rho_{\tau}^2} \]
Also
\[A_{\text{large}}=\sum_{i=1}^{n_{\text{large}}} a_{s_{i}}\] 
and all the circles in \(\mathcal{C}_{L}^*\) have radius \(\geq \rho_{\tau-1}\). For a fixed \(A_{\text{large}}\) and among all the possible candidate sets of \(\mathcal{C}_{L}^*\), \(n_{\text{large}}\) will be maximum when all the circles in \(S^*\) have radius \(\rho_{\tau-1}\). Hence we can bound \(n_{\text{large}}\) as \[n_{\text{large}} \leq \frac{A_{\text{large}}}{\pi \rho_{\tau-1}^2} \]
Also note that

\begin{align*}
A_{\text{small}} &> A_{\text{large}}\times \frac{\rho_{\tau-1}^2}{4}\\
\Rightarrow n_{\text{small}}\times \pi \rho_{\tau}^2 &\geq n_{\text{large}} \times \pi \rho_{\tau-1}^2 \times \frac{\rho_{\tau-1}^2}{30} \\
\Rightarrow \frac{n_{\text{small}}}{n_{\text{large}}} &\geq \frac{\rho_{\tau-1}^4}{30\rho_{\tau}^2}
\end{align*}

Also we know that 
\begin{align*}
\rho_{\tau} &= \frac{1}{2} \times \rho_{\tau-1}^2\eps 
\Rightarrow \frac{\rho_{\tau-1}^4}{4\rho_{\tau}^2} = \frac{1}{\eps^2} \\
\Rightarrow \frac{n_{\text{small}}}{n_{\text{large}}} &\geq \frac{1}{8{\eps}^2} 
\Rightarrow n_{\text{large}} \leq 8\eps^2 \times n_{\text{small}}
\end{align*}
As the number of large circles in \(\mathcal{C}_{L}^*\) is much less compared to the number of small circles in \(\mathcal{C}_{S}^*\) and all the circles have unit profit, we choose to pack only the small circles. We remove all the large circles from \(\mathcal{C}_{L}^*\) that we have already packed in the bin. We continue only with the small circles in \(\mathcal{C}_{S}^*\) and pack these using $\mathtt{greedyNFDR}$ for unit profit circles in our optimal solution. Our final packing \(P^*\) will have weight \((\frac{\pi}{4}(1-4\rho_{\tau})-\pi\rho_{\tau}^2)\lvert \mathcal{C}_{S}^* \rvert\geq (\frac{\pi}{4}(1-4\rho_{\tau})-\pi\rho_{\tau}^2)\frac{(1-\eps)}{(1+8\eps^2)}w(\OPT)\geq(\frac{\pi}{4}(1-4\frac{\eps}{8})-\pi(\frac{\eps}{8})^2)\frac{(1-\eps)}{(1+8\eps^2)}w(\OPT)\geq\frac{3}{5}w(\OPT)\) as \(\eps \leq 0.1\).

\begin{theorem}
	\label{thm:5/3-apx}
	There is a polynomial time 5/3-approximation algorithm for the unweighted $\mathsf{CSK}$ problem in two dimensions.
\end{theorem}
\begin{proof}
	First we divide our set $\mathcal{C}$ of circles into intervals as described in section \ref{partiotion section}. After the intervals are produced, we will consider each as of the first \(\lceil\frac{1}{\eps}\rceil\) intervals as medium separately and produce a packing as per the structure as defined above. This approach will produce a finite number of packings and we can compare all of them and then return the one with the highest weight.
	We do the following for all \(\tau\) such that \(1 \leq \tau \leq \frac{1}{\eps}\).
	For \(\tau = 1\), we essentially have only medium and small circles. We discard the medium circles and pack a set of small circles using $\mathtt{greedyNFDR}$ as described in lemma \ref{greedyNFDR_lemma}. Essentially this case corresponds to when \(\mathcal_{C}_{L}^*=\emptyset\) and as described in section \ref{Big_Circle_Pack} this strategy gives us $\frac{5}{3}$-approximation.
	For \(\tau \neq 1\), we have all three type of circles and we call the of large circles as \(\mathcal{C}_{L}\), the set of medium circles as \(\mathcal{C}_{M}\) and the set of small circles as \(\mathcal{C}_{S}\). First we note that all the circles in \(\mathcal{C}_L\) have radius \(\geq \rho_{\tau-1}\) and area \(\geq \pi \rho_{\tau-1}^2\). Hence in an optimal packing there can be at most \(\frac{1}{\pi \rho_{\tau-1}^2}\) in any feasible packing. We guess the large circles in \(\OPT\) in \(n^{\left(\frac{1}{\pi\rho_{\tau-1}^2}\right)^{\mathcal{O}(1)}} \)time which is essentially \(n^{({\frac{1}{\eps}})^{\mathcal{O}(1)}}\) time since \(\rho_{\tau-1}\) is bounded by \(\eps\). We can guess the \(\mathcal{C}_{L}^*\) by exhausting all the possible feasible packing. If we find any circle in \(\mathcal{C}_{L}^*\) has diameter \(\geq 1-2\rho_{\tau-1}\), then we simply discard this circle from the set \(\mathcal{C}\) and pack other circles in $\mathcal{C}$ using $\mathtt{greedyNFDR}$ as described in lemma \ref{greedyNFDR_lemma} as described in section \ref{Big_Circle_Pack} this strategy gives a packing weight \(w(\OPT)-1\). If no such circle exist then we follow the steps as described in section \ref{Big_Circle_Pack} and modify the packing to get the structure \(P^+\). We further use the ideas of section \ref{Small-Circle-Pack} to pack the small circles from the set \(\mathcal{C}_{S}\) where we choose greedily the circles with smaller radius from the set \(\mathcal{C}_S\) and arrive at the final packing \(P^*\). Note that choosing greedily the smaller circles from \(\mathcal{C}_{L}\) we either guess all or more circles from \(\mathcal{C}_{L}^*\) which essentially means our packing weight is at least as much as of the final packing \(P^*\) as in section \ref{Small-Circle-Pack}. We separately deal with the case when the number of small circles is much greater than the number of large circles. In this case we only pack the circles in \(\mathcal{C}_{S}\) using greedyNFDR in the unit bin. 
	Finally our algorithm returns the packing with the maximum weight among all the packings generated in polynomial time and from our analysis it follows that the best packing has profit at least \(\frac{3}{5}\cdot w(\OPT)\).
\end{proof}
\end{comment}